\section{Aplicação aos dados de treino}

\subsection{Caracterização do ruído}

Nas \(\valNNoise\) janelas de ruído puro, a média amostral foi
\(\valMuRuido\) e a variância \(\valVarRuido\) (desvio padrão
\(\valStdRuido\)). A autocorrelação média, normalizada, apresenta
\(\valFracAutocorrFora\%\) dos lags fora da faixa aproximada
\(\pm 2/\sqrt{N}\) esperada para ruído branco. A razão entre o máximo e o
mínimo da PSD estimada por Welch foi \(\valPsdRatio\).

Esses indícios levam a tratar o ruído como colorido
(\valNoiseColored). A covariância Toeplitz construída a partir da
autocorrelação empírica foi usada no BLUE colorido.

\begin{figure}[H]
\centering
\includegraphics[width=\textwidth]{01_ruido_autocorr_psd.png}
\caption{Autocorrelação média e PSD do ruído de treino.}
\label{fig:ruido}
\end{figure}

\subsection{Template do pulso}

O template concentra a maior parte da energia em baixas frequências: cerca de
\(\valEnergyLowBand\%\) da energia espectral está abaixo de \(100\,\mathrm{Hz}\).
A Figura~\ref{fig:template} compara o espectro normalizado do pulso com a PSD
do ruído. O conteúdo do pulso sobrepõe-se a uma região em que o ruído ainda é
relevante, o que ajuda a explicar a dificuldade da detecção em algumas janelas.

\begin{figure}[H]
\centering
\includegraphics[width=\textwidth]{02_template_espectro.png}
\caption{Forma temporal do template e comparação espectral com o ruído.}
\label{fig:template}
\end{figure}

\subsection{Filtro de Wiener}

O filtro de Wiener em frequência foi projetado com \(S_{xx}\) a partir do
template (escalado pela potência média das amplitudes do treino) e \(S_{rr}\)
estimada nas janelas de ruído puro, ambos na mesma convenção espectral. Como o
ruído tem potência relevante em toda a banda, o filtro resultante atenua
fortemente o sinal fora da faixa de baixas frequências em que o pulso concentra
energia. A Figura~\ref{fig:wiener} mostra exemplos antes e depois da filtragem:
a saída é bem mais suave, e o RMSE médio da forma, tomando a própria saída de
Wiener como estimativa do pulso, cai de \(\valRmseBrutoWave\) (sinal bruto)
para \(\valRmseWienerWave\). Há, porém, um efeito colateral: a saída é
atenuada em relação ao pulso original, o que prejudica a estimação direta da
amplitude a partir do máximo da janela filtrada, como se verá nas tabelas de
resultados.

\begin{figure}[H]
\centering
\includegraphics[width=\textwidth]{03_wiener_exemplos.png}
\caption{Exemplos de janelas com e sem pulso, antes e depois do Wiener.}
\label{fig:wiener}
\end{figure}

\subsection{Filtro casado}

A correlação com o template separa razoavelmente as estatísticas de janelas com
e sem pulso (Figura~\ref{fig:casado-sep}). O erro absoluto de posição tem
mediana \(\valNZeroErrMed\) amostras no treino, embora a média
(\(\valNZeroErrMean\)) seja maior por causa de outliers em que o máximo da
correlação não coincide com o início verdadeiro do pulso.

\begin{figure}[H]
\centering
\includegraphics[width=\textwidth]{04_casado_separacao.png}
\caption{Separação das estatísticas do casado e histograma do erro de \(n_0\).}
\label{fig:casado-sep}
\end{figure}

\begin{figure}[H]
\centering
\includegraphics[width=\textwidth]{05_casado_exemplo.png}
\caption{Exemplo de janela com pulso e saída do filtro casado.}
\label{fig:casado-ex}
\end{figure}

\subsection{BLUE da amplitude}

Com \(\hat{n}_0\) obtido pelo casado, o BLUE branco atingiu RMSE de amplitude
\(\valRmseBlueBranco\) no treino, enquanto o BLUE colorido reduziu o erro para
\(\valRmseBlueColorido\). Ambos superam o baseline ingênuo que usa o máximo de
\(y[n]\) (RMSE \(\valRmseNaiveA\)). A Figura~\ref{fig:blue} mostra a dispersão
entre amplitudes verdadeiras e estimadas.

\begin{figure}[H]
\centering
\includegraphics[width=\textwidth]{06_blue_scatter.png}
\caption{Amplitude estimada versus verdadeira (treino), BLUE branco e colorido.}
\label{fig:blue}
\end{figure}

\subsection{Limiar de detecção}

Usando apenas as janelas de ruído puro, o limiar do filtro casado foi definido
como o percentil 99 da estatística de detecção: \(\eta = \valThrCasado\). Por
construção, a taxa de falso alarme nesse mesmo conjunto é de aproximadamente
\(1\%\); a validação efetiva do limiar é feita em dados que não participaram
do ajuste, por meio da taxa de falso alarme empírica nas janelas sem pulso dos
conjuntos de treino e de teste (coluna ``FA emp.'' das
Tabelas~\ref{tab:treino} e~\ref{tab:teste}), que permanece próxima do alvo.

\begin{figure}[H]
\centering
\includegraphics[width=0.85\textwidth]{07_limiar.png}
\caption{Histograma da estatística sob \(H_0\) e limiar correspondente.}
\label{fig:limiar}
\end{figure}
